% Diagram and definition companion to ``Error corrections in view of
% algebraic topology''. Include in a document loading amsmath and amssymb.
% This is a fragment, not a standalone document; no TikZ is required.

\section*{Spaces, cellular algebra, and surface-code operators}

\subsection*{Quotient spaces and homotopy}

For a topological space $K$ and an equivalence relation $\sim$, the
projection $q:K\to K/{\sim}$ sends a point to its equivalence class.
The quotient topology is characterized by
\[
 U\subseteq K/{\sim}\text{ open}
 \quad\Longleftrightarrow\quad q^{-1}(U)\subseteq K\text{ open}.
\]
For a nonempty subspace $A\subseteq K$, the space $K/A$ identifies all
points of $A$ to one point. This topological quotient is different from
the quotient of chain modules introduced below.

For a based space $(K,k_0)$ and $n\geq1$, put
\[
 \pi_n(K,k_0):=[(S^n,*),(K,k_0)]_*.
\]
The brackets mean based maps modulo based homotopy: a homotopy
$F:S^n\times[0,1]\to K$ keeps $F(*,t)=k_0$ for every $t$.
For $n=1$, multiplication is loop concatenation. For path-connected
CW complexes,
\[
 H_1(K;\mathbb Z)\cong\pi_1(K,k_0)_{\mathrm{ab}},
 \qquad
 H_1(K;\mathbb F_2)\cong H_1(K;\mathbb Z)\otimes\mathbb F_2.
\]
Here abelianization is quotient by the commutator subgroup. Homotopy
classes of maps and homology classes of chains are distinct constructions.

\subsection*{Cellular modules and incidence maps}

Let $K$ be a finite oriented regular two-dimensional CW complex,
and let $R$ be a commutative ring with unit. Put
\[
 C_k:=C_k(K;R)=\bigoplus_{\sigma\in K_k}R\sigma,
 \qquad C^k:=\operatorname{Hom}_R(C_k,R).
\]
The set $K_k$ indexes the open $k$-cells. The dual basis satisfies
$\sigma^\vee(\tau)=\delta_{\sigma\tau}$; evaluation is the canonical pairing
$C^k\times C_k\to R$. Coordinate identifications require chosen bases.
Throughout this fragment $C_n(K;R)$ denotes cellular, rather than
singular, chains. The chosen cell orientations give an identification
\[
 C_n(K;R)\cong H_n(K^{(n)},K^{(n-1)};R),
 \qquad K^{(-1)}=\varnothing.
\]
Hatcher denotes this cellular differential by $d_n$; here it is
written $\partial_n$. The cellular homology theorem identifies the
homology of this complex with singular homology.
The chain complex and its dual are
\[
 C_2\xrightarrow{\partial_2}C_1\xrightarrow{\partial_1}C_0,
 \qquad
 C^0\xrightarrow{\delta^0}C^1\xrightarrow{\delta^1}C^2,
\]
where $\partial_1\partial_2=0$ and
$(\delta^k\phi)(c)=\phi(\partial_{k+1}c)$.
Thus $\delta^1\delta^0=0$.

For a face $f$, let $a_f:S^1\to K^{(1)}$ be its attaching map.
For an open edge $e$, let
$q_e:K^{(1)}\to K^{(1)}/(K^{(1)}\setminus e)\cong S^1$.
Using $\iota:\mathbb Z\to R$, $m\mapsto m1_R$, define
\[
 [f:e]:=\iota(\deg(q_e\circ a_f)),
 \qquad \partial_2f=\sum_e[f:e]e.
\]
If $e$ is oriented from $s(e)$ to $t(e)$, then
\[
 \partial_1e=t(e)-s(e),\qquad [e:v]:=v^\vee(\partial_1e),
 \qquad \delta^0v^\vee=\sum_e[e:v]e^\vee.
\]
The last formula takes a vertex \emph{cochain} as input.

\subsection*{Homology and exact sequences}

Define cycles, boundaries, and their quotient by
\[
 Z_k:=\ker\partial_k,\qquad B_k:=\operatorname{im}\partial_{k+1},
 \qquad H_k:=Z_k/B_k.
\]
Thus $z,z'\in Z_k$ have the same class exactly when
$z'-z=\partial_{k+1}b$ for some $b\in C_{k+1}$.
A complex requires $B_k\subseteq Z_k$; exactness at $C_k$ requires equality.
There are genuine short exact sequences
\[
 0\longrightarrow Z_k\xrightarrow{\ i\ }C_k
 \xrightarrow{\ \partial_k\ }B_{k-1}\longrightarrow0,
\]
\[
 0\longrightarrow B_k\xrightarrow{\ j\ }Z_k
 \xrightarrow{\ q\ }H_k\longrightarrow0.
\]
Here $i,j$ are inclusions, $q(z)=z+B_k$, and the first boundary map is
corestricted to its image. Dually,
\[
 H^k:=\ker\delta^k/\operatorname{im}\delta^{k-1}.
\]

A chain map $F_*$ from the cellular complex of $K$ to that of $L$
satisfies $\partial_k^L F_k=F_{k-1}\partial_k^K$:
\[
\begin{array}{ccccc}
 C_2(K;R)&\xrightarrow{\partial_2^K}&C_1(K;R)
 &\xrightarrow{\partial_1^K}&C_0(K;R)\\[3pt]
 {\scriptstyle F_2}\downarrow&&{\scriptstyle F_1}\downarrow
 &&{\scriptstyle F_0}\downarrow\\[3pt]
 C_2(L;R)&\xrightarrow{\partial_2^L}&C_1(L;R)
 &\xrightarrow{\partial_1^L}&C_0(L;R).
\end{array}
\]
It induces $H_k(F_*):H_k(K;R)\to H_k(L;R)$,
$[z]\mapsto[F_kz]$.
For a continuous cellular map $f:K\to L$, the induced cellular chain
map is denoted $f_\sharp$, while its induced homology map is $f_*$.

\subsection*{Relative chains and the connecting homomorphism}

Let $A\subseteq K$ be a CW subcomplex. Its chains form a subcomplex,
so the quotient carries a well-defined differential:
\[
 C_k(K,A;R):=C_k(K;R)/C_k(A;R),
\]
\[
 \partial_k^{\mathrm{rel}}:C_k(K,A;R)\longrightarrow C_{k-1}(K,A;R),
 \qquad \overline c\longmapsto\overline{\partial_kc}.
\]
Here $\overline c=c+C_k(A;R)$. A relative cycle satisfies
$\partial_kc\in C_{k-1}(A;R)$, rather than necessarily $\partial_kc=0$.
The relative homology and cochain modules are
\[
 H_k(K,A;R):=\ker\partial_k^{\mathrm{rel}}
 /\operatorname{im}\partial_{k+1}^{\mathrm{rel}},
\]
\[
 C^k(K,A;R):=\operatorname{Hom}_R(C_k(K,A;R),R)
 \cong\{\phi\in C^k:\phi|_{C_k(A;R)}=0\}.
\]
For a nonempty CW subcomplex,
$H_k(K,A;R)\cong\widetilde H_k(K/A;R)$.
The degreewise short exact sequence of chain complexes is
\[
 0\longrightarrow C_*(A;R)\xrightarrow{\ i\ }C_*(K;R)
 \xrightarrow{\ j\ }C_*(K,A;R)\longrightarrow0.
\]
Here $i$ is inclusion and $j$ is the quotient projection.
It induces the long exact sequence (coefficients $R$ suppressed)
\[
 \cdots\longrightarrow H_k(A)\xrightarrow{i_*}H_k(K)
 \xrightarrow{j_*}H_k(K,A)\xrightarrow{\partial}
 H_{k-1}(A)\xrightarrow{i_*}H_{k-1}(K)\longrightarrow\cdots.
\]
Following Hatcher, the connecting homomorphism is also denoted
$\partial$:
\[
 \partial:H_k(K,A;R)\longrightarrow H_{k-1}(A;R),
 \qquad [\overline c]\longmapsto[\partial_kc]_A.
\]
Indeed, lift a relative cycle to $c\in C_k(K;R)$; its boundary lies in
$A$ and is a cycle there. Replacing the lift and representative by
$c'=c+a+\partial_{k+1}b$, with $a\in C_k(A;R)$, changes this boundary
by $\partial_ka$, hence not its homology class in $A$.

The maps $\partial_k^{\mathrm{rel}}$ and $\partial$ have different domains
and codomains: the first acts on relative chains, the second on relative
homology. Neither is the boundary map of a relative homotopy sequence.

\subsection*{From cellular maps to operator products}

Choose additive-group homomorphisms
\[
 \rho:(C_1,+)\longrightarrow\operatorname{GL}(\mathcal H),
 \qquad \Pi:(C^1,+)\longrightarrow\operatorname{GL}(\mathcal H).
\]
These are not $R$-linear maps into $\operatorname{GL}(\mathcal H)$.
The typed compositions are
\[
 C_2\xrightarrow{\partial_2}C_1\xrightarrow{\rho}\operatorname{GL}(\mathcal H),
 \qquad f\longmapsto\partial_2f\longmapsto B_f,
\]
\[
 C^0\xrightarrow{\delta^0}C^1\xrightarrow{\Pi}\operatorname{GL}(\mathcal H),
 \qquad v^\vee\longmapsto\delta^0v^\vee\longmapsto A_v.
\]
With $W_e(r):=\rho(re)$ and $U_e(r):=\Pi(re^\vee)$,
\[
 B_f:=\rho(\partial_2f)=\prod_e W_e([f:e]),
 \qquad A_v:=\Pi(\delta^0v^\vee)=\prod_e U_e([e:v]).
\]
These compositions are neither chain complexes nor exact sequences.
Locality, unitarity, and commutation between the two families require
representation data.

\subsection*{Binary realization, syndromes, and logical classes}

Over $\mathbb F_2$, let $D_*=C_*(K;\mathbb F_2)$ for a closed surface,
or $D_*=C_*(K,A;\mathbb F_2)$ for a patch with rough boundary $A$.
Write $d_k$ for its differential and $d^k$ for the dual. Assign a qubit
to every cellular basis edge of $D_1$ and put $n=\dim D_1$. Then
\[
 D^k:=\operatorname{Hom}_{\mathbb F_2}(D_k,\mathbb F_2),
 \qquad (d^kx)(c):=x(d_{k+1}c),
 \quad x\in D^k,\ c\in D_{k+1}.
\]
The binary operator realization is
\[
 \mathcal H=(\mathbb C^2)^{\otimes n},
 \qquad \rho(z)=\bigotimes_e Z^{z_e},
 \qquad \Pi(x)=\bigotimes_e X^{x_e}.
\]
For $x\in D^1$ and $z\in D_1$, the Pauli relation is
\[
 \Pi(x)\rho(z)=(-1)^{x(z)}\rho(z)\Pi(x).
\]
Consequently, for $B_f=\rho(d_2f)$ and $A_v=\Pi(d^0v^\vee)$,
\[
 (d^0v^\vee)(d_2f)=v^\vee(d_1d_2f)=0
 \quad\Longrightarrow\quad A_vB_f=B_fA_v.
\]
The syndromes and pure-type logical Pauli classes are
\[
 s_Z=d_1z\in D_0,\qquad s_X=d^1x\in D^2,
\]
\[
 \mathcal L_Z\cong\ker d_1/\operatorname{im}d_2=H_1(D_*),
 \qquad \mathcal L_X\cong\ker d^1/\operatorname{im}d^0=H^1(D_*).
\]
A correction $\widehat z$ with $d_1\widehat z=d_1z$ succeeds exactly when
$[z+\widehat z]=0$ in $H_1(D_*)$. In cellular coordinates,
\[
 H_X=[d_1],\qquad H_Z=[d_2]^{\mathsf T},\qquad
 H_XH_Z^{\mathsf T}=0.
\]

For a disk $K$ with two disjoint contractible rough arcs
$A=A_0\sqcup A_1$, the pair sequence gives
\[
 0\longrightarrow H_1(K,A;\mathbb F_2)\xrightarrow{\partial}
 \mathbb F_2^2\xrightarrow{\ (a,b)\mapsto a+b\ }\mathbb F_2
 \longrightarrow0.
\]
Hence $H_1(K,A;\mathbb F_2)\cong\mathbb F_2$: a path joining the two
rough arcs represents the logical $Z$ class. Complementary smooth
boundaries describe the dual paths for logical $X$, with the appropriate
dual cellulation and consistent boundary assignments.

\subsection*{Invariant subspaces and correctable errors}

Let $S=\langle A_v,B_f\rangle\subset U(\mathcal H)$ be the finite
abelian Pauli stabilizer group, with $-I\notin S$. Its invariant
subspace and orthogonal projector are
\[
 \mathcal C=\mathcal H^S
 :=\{\psi\in\mathcal H:g\psi=\psi\text{ for all }g\in S\},
 \qquad
 P_{\mathcal C}=\frac1{|S|}\sum_{g\in S}g.
\]
Indeed, closure under inverses gives $P_{\mathcal C}^{\dagger}=P_{\mathcal C}$;
counting products gives $P_{\mathcal C}^2=P_{\mathcal C}$.
Multiplication by a group element permutes the sum, so
$\operatorname{im}P_{\mathcal C}=\mathcal H^S$.
The Pauli trace identity yields $\dim\mathcal C=2^n/|S|>0$.

More generally, let $G$ be a finite group, let
$U:G\to U(\mathcal H)$ be a unitary representation, and let
$\chi:G\to\mathbb C$ be an irreducible character of dimension
$d_\chi=\chi(e)$. The orthogonal isotypic projector is
\[
 P_\chi=\frac{d_\chi}{|G|}
 \sum_{g\in G}\overline{\chi(g)}\,U(g).
\]
Its image is the sum of the irreducible subrepresentations with
character $\chi$. For $G=S$, with $U$ the inclusion and $\chi=1$,
this reduces to the stabilizer code projector. Thus the cellular
maps construct the checks, while group averaging selects their
invariant quantum states.

An isotypic subspace alone need not correct the specified noise.
For a noise channel with Kraus operators
$E_a\in\operatorname{End}(\mathcal H)$, exact recovery on
$\mathcal C$ is equivalent to the Knill--Laflamme condition
\[
 P_{\mathcal C}E_a^{\dagger}E_bP_{\mathcal C}
 =c_{ab}P_{\mathcal C}\qquad\text{for every }a,b,
\]
where $c_{ab}\in\mathbb C$ is independent of the encoded state.
See E. Knill and R. Laflamme, \emph{A Theory of Quantum Error-Correcting
Codes}, Physical Review A \textbf{55}, 900 (1997),
arXiv:quant-ph/9604034.

\subsection*{A worked planar patch}

Let $K=[0,2]\times[0,2]$ be divided into four unit-square faces.
Let $A=A_L\sqcup A_R$ consist of its left and right sides, including
their vertices and edges. These are the rough boundary arcs; the
complementary top and bottom arcs are smooth, with corner points in $A$.
We use the convention that $Z$ paths
may terminate on rough boundaries, whereas dual $X$ paths may
terminate on smooth boundaries.

All coefficients are now in $\mathbb F_2$. Abbreviate
\[
 C_n=C_n(K,A;\mathbb F_2),\qquad
 C^n=\operatorname{Hom}_{\mathbb F_2}(C_n,\mathbb F_2),
\]
and denote the induced relative differential by
$\partial_n:C_n\to C_{n-1}$. A bar denotes a relative chain class,
for example $\bar e=e+C_1(A;\mathbb F_2)$; brackets $[z]$ denote a
homology class only when $z$ is a relative cycle.

Index rows from top to bottom by $i=0,1,2$. Write $h_{iL},h_{iR}$
for the horizontal edges joining the left, middle, and right columns.
The middle vertical edges are $m_0,m_1$, and the middle-column
vertices are $v_0,v_1,v_2$. The faces between rows $i$ and $i+1$ are
$f_{iL},f_{iR}$, for $i=0,1$. Choose ordered bases
\[
 \mathcal B_2=(\bar f_{0L},\bar f_{0R},\bar f_{1L},\bar f_{1R}),
 \qquad \mathcal B_0=(\bar v_0,\bar v_1,\bar v_2),
\]
\[
 \mathcal B_1=(\bar h_{0L},\bar h_{0R},\bar h_{1L},\bar h_{1R},
 \bar h_{2L},\bar h_{2R},\bar m_0,\bar m_1).
\]
The four vertical boundary edges and their six vertices vanish in
the relative quotient. Hence
\[
 C_2\cong\mathbb F_2^4,\qquad C_1\cong\mathbb F_2^8,
 \qquad C_0\cong\mathbb F_2^3.
\]
There are eight data qubits, one for each basis edge of $C_1$.

The relative boundary formulas are
\[
 \partial_1\bar h_{iL}=\partial_1\bar h_{iR}=\bar v_i
 \quad(0\leq i\leq2),\qquad
 \partial_1\bar m_i=\bar v_i+\bar v_{i+1}\quad(0\leq i\leq1),
\]
\[
 \partial_2\bar f_{iL}=\bar h_{iL}+\bar h_{i+1,L}+\bar m_i,
 \qquad
 \partial_2\bar f_{iR}=\bar h_{iR}+\bar h_{i+1,R}+\bar m_i
 \quad(0\leq i\leq1).
\]
With coordinate columns in the stated bases,
\[
 [\partial_1]_{\mathcal B_0\leftarrow\mathcal B_1}
 =\begin{pmatrix}
 1&1&0&0&0&0&1&0\\
 0&0&1&1&0&0&1&1\\
 0&0&0&0&1&1&0&1
 \end{pmatrix},
\]
\[
 [\partial_2]_{\mathcal B_1\leftarrow\mathcal B_2}
 =\begin{pmatrix}
 1&0&0&0\\
 0&1&0&0\\
 1&0&1&0\\
 0&1&0&1\\
 0&0&1&0\\
 0&0&0&1\\
 1&1&0&0\\
 0&0&1&1
 \end{pmatrix}.
\]
Direct multiplication and row reduction give
\[
 [\partial_1][\partial_2]=0_{3\times4},\qquad
 \operatorname{rank}[\partial_1]=3,\qquad
 \operatorname{rank}[\partial_2]=4.
\]
Therefore $\dim H_1(K,A;\mathbb F_2)=(8-3)-4=1$. The three
independent star checks and four independent plaquette checks
encode one qubit. For instance,
\[
 A_{v_1}=\Pi(\delta^0\bar v_1^\vee)
 =X_{h_{1L}}X_{h_{1R}}X_{m_0}X_{m_1},
 \qquad
 B_{f_{0L}}=\rho(\partial_2\bar f_{0L})
 =Z_{h_{0L}}Z_{h_{1L}}Z_{m_0}.
\]

\paragraph{A shared syndrome and two recoveries.}
Take the actual $Z$-error label $z=\bar h_{1L}$. The proposed
recovery $\widehat z=\bar h_{1R}$ has the same syndrome:
\[
 s_Z=\partial_1z=\partial_1\widehat z=\bar v_1,
 \qquad [s_Z]_{\mathcal B_0}=\begin{pmatrix}0\\1\\0\end{pmatrix}.
\]
Its residual is the relative cycle
\[
 \ell_Z=z+\widehat z=\bar h_{1L}+\bar h_{1R},
 \qquad \partial_1\ell_Z=0.
\]
Since $K$ is contractible and $A$ has two contractible components,
the exact sequence of the pair reduces to
\[
 0\longrightarrow H_1(K,A;\mathbb F_2)
 \xrightarrow{\partial}H_0(A;\mathbb F_2)
 \xrightarrow{i_*}H_0(K;\mathbb F_2)\longrightarrow0.
\]
Here $i_*(a,b)=a+b$ in component coordinates. The ordinary lift of
$\ell_Z$ has one endpoint on each rough arc, so
\[
 \partial[\ell_Z]=[a_L]+[a_R]\ne0,
 \qquad a_L\in A_L,\quad a_R\in A_R.
\]
Thus $[\ell_Z]\ne0$: this recovery removes the syndrome but leaves
a logical error. The connecting map $\partial$ acts on homology
classes; the relative chain differential $\partial_1$ acts on chains.

A successful recovery for the same error is
\[
 \widehat z_{\mathrm{good}}=\bar h_{0L}+\bar m_0,
 \qquad \partial_1\widehat z_{\mathrm{good}}=\bar v_1,
 \qquad z+\widehat z_{\mathrm{good}}=\partial_2\bar f_{0L}.
\]
Its residual operator is the stabilizer $B_{f_{0L}}$, which acts as
the identity on the code space. The recovery need not retrace the
error; its residual must be a relative boundary.

\paragraph{Homological deformation and the conjugate logical.}
The top representative $\ell'_Z=\bar h_{0L}+\bar h_{0R}$ satisfies
\[
 \ell'_Z+\ell_Z=\partial_2(\bar f_{0L}+\bar f_{0R}),
 \qquad [\ell'_Z]=[\ell_Z]\in H_1(K,A;\mathbb F_2).
\]
For a conjugate logical take the cochain
\[
 x=\bar h_{0L}^{\vee}+\bar h_{1L}^{\vee}+\bar h_{2L}^{\vee}
 \in C^1(K,A;\mathbb F_2).
\]
It evaluates to zero on every face boundary and to one on $\ell_Z$:
\[
 \delta^1x=x\circ\partial_2=0,\qquad x(\ell_Z)=1.
\]
Hence both classes are nonzero, and the logical operators
\[
 \overline X=\Pi(x),\qquad \overline Z=\rho(\ell_Z)
 \quad\text{satisfy}\quad
 \overline X\,\overline Z=-\overline Z\,\overline X.
\]
The dual path depicting $x$ joins the smooth boundaries and crosses
the primal representative $\ell_Z$ once. Its algebraic label remains
a cochain, not a primal edge chain. This small patch has
$d_Z=2$, $d_X=3$, and parameters $[[8,1,2]]$.

\paragraph{Reference for notation.}
Allen Hatcher, \emph{Algebraic Topology}, Chapter~2, pp.~115--118
(relative chains and the connecting boundary), and Theorem~2.35,
p.~139 (cellular homology); Chapter~3, pp.~197--198 (cochains and
coboundaries).
